\uuid{5Q8c}
\titre{ Étude d'une application linéaire }

\niveau{L2}
\module{ Applications linéaires }
\niveau{} 				%L1, L2, L3, MPSI, MP, PCSI, PC, PSI...
\module{ Espaces vectoriels } 	%Analyse, Algèbre...
\chapitre{Réduction d'endomorphisme}   			%Continuité, Groupes, Fonctions de plusieurs variables...
\sousChapitre{Diagonalisation}			%Optimisation, Diagonalisation d'une matrice, Calcul de dérivées partielles...

\theme{}				%Fonction de répartition, Division euclidienne de polynômes, ...
\auteur{Quentin Liard}
\datecreate{ 2026-02-19}
\organisation{AMSCC}			%AMSCC, Exo7, ...
\difficulte{1}			%1, 2, 3, 4 ou 5






\contenu{

\texte{
On considère l’application linéaire $f:\mathbb{R}^3\rightarrow\mathbb{R}^3$ définie par

\[
f(x,y,z)=(2x+y,x+2y,z).
\]
}

\begin{enumerate}

\item \question{Déterminer la matrice de $f$ dans la base canonique.}

\item \question{Calculer le polynôme caractéristique de cette matrice.}

\item \question{Déterminer les valeurs propres et les sous-espaces propres.}

\item \question{Montrer que $f$ est diagonalisable.}

\item \question{Déterminer une base de vecteurs propres de $\mathbb{R}^3$.}

\end{enumerate}

}